\subsection{SUBGROUPS OF A TOPOLOGICAL GROUP}

Let G be a topological group and let H be a subgroup of G. By (GT'), the topology induced on H by the topology of G is compatible with the group structure of H. The structure of a topological group thus defined on H is said to be induced by that of G. Whenever we consider a subgroup H of G as a topological group, it is always this induced structure which is under consideration, unless the contrary is expressly stated.

PROPOSITION 1. The closure \(\overline{H}\) of a subgroup H of a topological group G is a subgroup of G. If H is a normal subgroup of G, then so is \(\overline{H}\) .

If \(a, b \in \overline{\mathbf{H}}\) then \(ab^{-1} \in \overline{\mathbf{H}}\) , because the mapping \((x, y) \to xy^{-1}\) is continuous on \(\mathbf{G} \times \mathbf{G}\) and transforms \(\mathbf{H} \times \mathbf{H}\) into \(\mathbf{H}\) (Chapter I, § 2, no. 1, Theorem 1). In the same way, the continuity of the mapping \(x \to axa^{-1}\) shows that if \(\mathbf{H}\) is normal then \(\overline{\mathbf{H}}\) is normal.

In particular, the closure N of the set \(\{e\}\) consisting only of the identity element of G, is a normal subgroup of G; and \(N=\{e\}\) if and only if G is Hausdorff (§ 1, no. 2, Proposition 2).

PROPOSITION 2. If G is a Hausdorff topological group, then the closure of a commutative subgroup of G is a commutative subgroup of G.

By reason of Proposition 1, we may limit ourselves to the case where H is dense in G. The continuous functions xy and yx are equal on H × H, and are therefore equal on G × G, by virtue of the principle of extension of identities (Chapter I, § 8, no. 1, Proposition 2, Corollary 1).

PROPOSITION 3. Let G be a Hausdorff topological group and let M be any subset of G. Then the set M' of elements of G which commute with each element of M is a closed subgroup of G. In particular, the centre of G is closed in G.

For \(\mathbf{M}'\) is the intersection of the sets \(\mathbf{F}_m(m \in \mathbf{M})\) , where \(\mathbf{F}_m\) is the set of all \(x \in \mathbf{G}\) such that \(xm = mx\) ; and \(\mathbf{F}_m\) is closed (Chapter I, § 8, no. 1, Proposition 2).

PROPOSITION 4. Let G be a topological group and let H be a subgroup of G which is locally closed at one point of H (Chapter I, § 3, no. 3, Definition 2). Then H is closed in G.

By translation, H is locally closed at each of its points, i.e.~H is locally closed in G. Let V be a symmetric open neighbourhood of e in G such that \(V \cap H\) is closed in V. If \(x \in \overline{H}\) , then xV meets H; if \(y \in xV \cap H\) , we have \(x \in yV\) , and \(y(V \cap H) = (yV) \cap H\) is closed in yV. But x lies in the closure of \((yV) \cap H\) , hence \(x \in H\) .

COROLLARY. A subgroup of a topological group is open if and only if it has an interior point. Every open subgroup is closed.

If a subgroup H has an interior point, then by translation all the points of H are interior and so H is open. The second assertion of the corollary is a particular case of Proposition 4.

PROPOSITION 5. A subgroup H of a topological group G is discrete if and only if H has an isolated point. Every discrete subgroup of a Hausdorff group is closed.

If H is discrete, every point of H is isolated. Conversely if H has an isolated point, then by translation every point of H is isolated and therefore H is discrete. If H is discrete and G is Hausdorff, then there is a neighbourhood V of e such that \(V \cap H = \{e\}\) ; \(\{e\}\) is closed in G and therefore a fortiori in V, so that H is locally closed at e. Hence H is closed in G by Proposition 4.

Remark. Let H be any subgroup of a topological group G. For each \(x \in \overline{H}\) we have \(\overline{xH} = x\) . \(\overline{H} = \overline{H}\) , since translation by x is a homeomorphism of G onto G. In other words, for each \(x \in \overline{H}\) , xH is dense in \(\overline{H}\) . It follows that if H is not closed, then \(\overline{H} \cap \mathrm{C}H\) is dense in \(\overline{H}\) .

